% !TEX root = ../Thesis.tex
\chapter{Methodology} \label{cha:method}

% \section{Choice of Research Strategy}



% \
% This thesis adopts Design Science Research (DSR) as its overall research framework. Design Science Research (DSR) is a research framework concerned with the design and evaluation of artefacts intended to address practical problems \cite{johannesson2021,hevner2004design}. In information systems research, DSR is particularly relevant when the aim is not only to solve a real-world problem, but also to generate design knowledge through the artefact development process \cite{gregor2013}. 

% This research fits that orientation well. Rather than only analyzing the limitations of the existing SAP-based TMC implementation at Volvo CE, the thesis aims to design and implement a web-based prototype that improves transparency, usability, and traceability in the preparation and adjustment stages. The study therefore focuses on creating an artefact that addresses a concrete industrial problem and examining its usefulness in practice.

% To structure the research process, the thesis follows the five core DSR activities proposed by Johannesson and Perjons: explicate the problem, define requirements, design and develop the artefact, demonstrate the artefact, and evaluate the artefact \cite{johannesson2021,hevner2004design}. In this thesis, these activities provide a systematic structure for moving from problem understanding to requirement definition, prototype development, and formative evaluation in the Volvo CE context.

% The five activities applied in this thesis are summarized as follows:

% \begin{itemize}
%     \item \textbf{Explicate the problem}: identify and clarify the limitations of the current TMC workflow and the needs of potential users.
%     \item \textbf{Define requirements}: derive functional and non-functional requirements for the prototype based on workflow analysis, documentation, and stakeholder needs.
%     \item \textbf{Design and develop the artefact}: design and implement a web-based frontend-backend prototype for data upload, workflow navigation, calculation-layer visualization, and interactive table analysis.
%     \item \textbf{Demonstrate the artefact}: apply the prototype to representative TMC-related data and workflow scenarios to show its practical usefulness.
%     \item \textbf{Evaluate the artefact}: assess whether the prototype improves visibility, usability, and interpretability of the preparation and adjustment stages.
% \end{itemize}

% Using DSR makes it possible to connect practical system development with academic research in a structured way. It provides a methodological basis for both building the prototype and reflecting on its contribution as a reusable approach for improving workflow transparency in enterprise data-processing systems.


% \subsection{Choice of Research Framework}

% Design Science Research (DSR) is a research strategy concerned with the design and evaluation of artefacts intended to address practical problems in real-world contexts \cite{johannesson2021,hevner2004design}. This thesis adopts DSR as its overall research framework, as the study aims not only to examine limitations in the existing SAP-based Total Market Calculation (TMC) workflow at Volvo CE, but also to design and evaluate a prototype that improves its transparency, usability, and traceability.

% DSR is particularly suitable for this thesis because the main contribution is an artefact rather than a purely descriptive account of the current workflow. While approaches such as case study are useful for understanding organizational settings and existing practices, they do not by themselves provide a structured process for developing and evaluating a technical solution. In contrast, DSR offers a systematic way to move from problem identification to requirement definition, artefact design, demonstration, and evaluation \cite{johannesson2021,hevner2004design}.

% In this thesis, the artefact is a web-based prototype that reconstructs and visualizes selected preparation and adjustment stages of the TMC workflow outside the original SAP environment. The purpose of the prototype is to make key transformations, data dependencies, and intermediate outputs easier to inspect and understand. In this way, the artefact addresses a practical industrial problem while also providing knowledge about how workflow transparency can be improved in enterprise data-processing systems.

% To structure the research process, the thesis follows the five core DSR activities proposed by Johannesson and Perjons: explicate the problem, define requirements, design and develop the artefact, demonstrate the artefact, and evaluate the artefact \cite{johannesson2021,hevner2004design}. In the present study, these activities provide a clear methodological structure for moving from understanding the limitations of the current TMC process to developing and assessing a prototype in the Volvo CE context.


% \section{Research Strategy}

Design Science Research (DSR) was selected because this thesis was not limited to describing the existing SAP-based TMC implementation. The main task was to build an artefact, apply it to realistic case material, and evaluate whether it made selected process stages easier to inspect \cite{johannesson2021,hevner2004design}. In DSR terms, the artefact is the proposed inspection-oriented solution, while the implemented prototype is its concrete instantiation in this thesis. This orientation fits the study because the intended contribution is partly practical and partly explanatory: the work needed to solve a real workflow-inspection problem, but also to show what kind of design choices made that solution useful in context.

The work followed the five DSR activities described by Johannesson and Perjons \cite{johannesson2021}, but not as a rigid linear sequence. Problem understanding, requirement clarification, implementation, and evaluation informed each other during the project. In practice, this mattered because several parts of the workflow could not be understood from documentation alone. As prototype outputs became available, they also revealed ambiguities that had to be discussed again with stakeholders before later stages could be implemented with confidence.

This thesis combines Design Science Research (DSR) with a case study. \textbf{Alternative research methods.} A stand-alone \textbf{case study} would have been appropriate for describing the organizational setting and the existing workflow \cite{yin2018}, but this thesis also required artefact design and evaluation. \textbf{Action research} was less suitable because the project did not aim to intervene directly in Volvo CE's live production process \cite{susman1978}. More measurement-oriented approaches, such as \textbf{surveys} or \textbf{controlled experiments}, were also less suitable because the work depended on a small group of knowledgeable stakeholders and on detailed interpretation of workflow logic.

The broader design problem concerns transparency, traceability, and usability in ERP-based enterprise data processing, but it was studied through one concrete setting: the TMC process at Volvo CE. Within the overall DSR design, the Volvo CE setting functioned as a case study used to understand the workflow problem in context, guide artefact design, and assess the implemented prototype in a real organizational environment.



% \section{Application of the Design Science Research Framework}
\section{Research Method}


% This thesis used a combination of qualitative and technical data collection methods across the five DSR activities. DSR supports the development and evaluation of artefacts in response to practical problems, while allowing different empirical methods to be used depending on the purpose of each activity \cite{johannesson2014introduction,hevner2004design}. Since the study was conducted in a specialised organisational context, document review, workflow inspection, stakeholder discussions, and prototype-based feedback were used as the main data collection methods. These methods are suitable for small-scale qualitative research where the aim is to obtain contextual and practice-oriented understanding rather than broad statistical generalisation \cite{denscombe2017good}.

% In the problem explication stage, data were collected through review of TMC-related documentation, inspection of the existing SAP-based implementation, and informal discussions with stakeholders in Volvo CE Global Business Intelligence. These sources were used to understand the limitations of the current workflow, particularly the lack of transparency, traceability, and inspectability in intermediate calculation steps. In the requirements definition stage, additional data were collected from workflow documentation, source data structures, matrix files, example outputs, and stakeholder input. This helped translate the identified problem into functional, structural, and environmental requirements for the artefact.

% During the design and development stage, data collection focused on technical material needed to reconstruct selected parts of the TMC workflow. This included CSV-based source files, matrix tables, mapping rules, and selected SAP-related workflow descriptions. Clarifications were also collected through stakeholder discussions when calculation rules or data dependencies required further interpretation. In the demonstration stage, the artefact was applied to a realistic TMC workflow scenario using 2024 source data and matrix files. The generated preparation, market view, adjustment, machine line split, and resplit outputs were used to demonstrate how the prototype could make intermediate workflow layers more inspectable. Finally, in the evaluation stage, data were collected through prototype-based stakeholder feedback, semi-structured discussion, output inspection, and requirement-based review.

% Alternative data collection methods were also considered. Large-scale surveys, focus groups, formal requirements workshops, direct SAP backend extraction, production-system integration, and full operational pilot deployment could have been used. However, these methods were less suitable for the scope of this thesis. The TMC workflow is a specialised internal process involving a limited number of knowledgeable users, and the artefact was developed as a research prototype outside the production SAP environment. Therefore, qualitative and technical data collection methods were considered more appropriate than broad survey-based or production-level methods.

% \subsection{Data Collection Methods}

This thesis used three main sources of empirical material: \textbf{document analysis}, \textbf{stakeholder clarification and feedback sessions}, and \textbf{prototype-based material}. To understand the TMC process, it was necessary to review workflow descriptions, source files, and matrix structures, but some rule interpretations still depended on business context.

Document analysis was therefore used to examine process materials, workflow descriptions, source data structures, and matrix files \cite{bowen2009document}. The document set was limited to materials directly relevant to the studied 2024 TMC workflow and the implemented stages. Outdated versions, unrelated SAP reports, and partial notes that could not be connected to the selected workflow scope were not treated as primary analytical material. This helped keep the analysis close to the actual workflow represented in the artefact, rather than drifting into broader SAP material that was not operationally relevant for the case. When documents conflicted or left rule details unclear, stakeholder clarification and output comparison were used to resolve the interpretation.

Recurring weekly review meetings and related clarification discussions were used to clarify business rules, expected outputs, and gaps in the documentation. These interactions were not conducted as a formal interview study. Instead, they functioned as case-based clarification and feedback sessions within the broader DSR and case study design. To avoid relying only on individual interpretations, unclear points were checked against documents, source structures, matrix definitions, and generated outputs wherever possible. Prototype-based material was then collected during design, demonstration, and evaluation in the form of uploaded inputs, generated outputs, result comparisons, and stakeholder comments \cite{runeson2009guidelines}.





% \subsection{Data Analysis Method}

The collected material was analysed through a combination of \textbf{qualitative interpretation} and \textbf{technical inspection}. In practice, this meant breaking the workflow into stages, interpreting rules from documents and stakeholder explanations, checking how those rules appeared in prototype outputs, and comparing implemented results with expected behaviour \cite{johannesson2021,hevner2004design}. The study was not conducted as a formal line-by-line content analysis. Instead, the main unit of analysis was the \textit{workflow stage} together with its associated inputs, rules, transformations, and generated outputs.

The analytical categories were mainly deductive. Transparency, traceability, usability, validation support, and requirement fulfilment were derived from the research problem, related work, and the DSR purpose of the study. At the same time, the analysis remained open to case-specific issues such as unclear rule ownership, reference-data version differences, and stage-specific inspection needs. Accordingly, the analysis should be understood as structured workflow analysis in a DSR case study rather than as an inductive code-to-theme thematic analysis.

During problem explication and requirements definition, the analysis focused on understanding where business users lost visibility in the existing process and which functions would be most useful in an inspection-oriented prototype. During design and development, the analysis became more concrete: source files, mapping tables, matrix rules, and workflow descriptions were examined to determine how selected SAP-based logic could be reconstructed as explicit processing steps. During demonstration and evaluation, the generated outputs were inspected stage by stage and compared with expected workflow behaviour and selected reference outputs. Stakeholder feedback was reviewed qualitatively, especially comments about whether the prototype made it easier to understand a stage, inspect an intermediate result, or discuss an unexpected value.







\section{Application of the Research Method}

This section summarizes how the selected data collection and analysis methods were applied across the five DSR activities. Detailed artefact design, implementation, demonstration, and evaluation are presented in later chapters.

\subsection{Explicate the Problem}

The first DSR activity was to explicate the problem. This was done by reviewing workflow descriptions, source structures, matrix files, existing SAP-based outputs, and stakeholder explanations in order to understand where business users lost visibility in the current TMC process. The goal was not only to describe how the workflow runs, but to identify where input dependencies, transformation logic, and intermediate results became difficult to inspect.

The collected material was analysed through qualitative interpretation of the workflow steps and technical inspection of how data and rules were distributed across the existing implementation. This made it possible to frame the problem more precisely as one of transparency, traceability, and usability in a staged enterprise data-processing workflow rather than as a single isolated system defect.

\subsection{Define Requirements}

The second DSR activity was to define requirements for the artefact. Workflow documentation, source data structures, matrix files, existing output examples, and stakeholder input were used to clarify what the prototype needed to support. The purpose was to translate the identified workflow problems into concrete requirements for a web-based artefact.

The collected material was analysed through requirements analysis. The requirements were grouped into functional and performance requirements. Functional requirements described capabilities such as file upload, backend processing, structured storage, intermediate output inspection, filtering, validation support, and export. Performance requirements addressed prototype-level execution time, stability, and reliability during demonstration and evaluation.

The resulting requirements guided the design and development of the prototype and provided the basis for the requirement-based evaluation in Chapter 6, where the artefact is assessed in terms of functional coverage, transparency, traceability, usability, validation support, and prototype-level performance.

















% The third activity of the DSR framework is to design and develop the artefact. This activity takes the outlined artefact and the elicited requirements as input and transforms them into a concrete implementation. According to Johannesson and Perjons, this activity can be understood through four sub-activities: imagine and brainstorm, assess and select, sketch and build, and justify and reflect \cite{johannesson2021,hevner2004design}. In this thesis, the activity was carried out through a combination of iterative prototyping, workflow analysis, and discussions with relevant stakeholders in the Volvo CE context.

% The first sub-activity, \textit{imagine and brainstorm}, involved exploring how the TMC workflow could be represented in a more transparent and user-friendly way outside the current SAP-based environment. Early ideas focused on how source data upload, workflow navigation, intermediate layer inspection, and calculation-related views could be presented through a web-based interface. These ideas were developed through technical analysis of the workflow and informal discussions with stakeholders involved in business intelligence and market-related analysis.

% The second sub-activity, \textit{assess and select}, involved evaluating alternative design directions and selecting the ones most suitable for the scope of the thesis. A web-based prototype was chosen rather than a script-based or spreadsheet-based solution, as it provided better support for structured navigation, interactive data inspection, and future internal use within Volvo CE. In addition, a frontend-backend architecture was selected in order to separate user interaction from backend processing and data handling. The implementation scope was defined up to the machine line split stage, while the preparation and adjustment stages were given particular emphasis in the thesis.

% The third sub-activity, \textit{sketch and build}, involved translating the selected design into a working prototype. The artefact was implemented as a web-based system with dedicated views for authentication, homepage navigation, source data upload, and workflow layer inspection. Uploaded CSV-based files were processed through backend services and stored in a local relational database, allowing selected parts of the workflow logic to be re-expressed through Python- and SQL-based processing in a more inspectable form than in the original SAP environment. The prototype was developed iteratively, and different parts of the system were refined as the understanding of the workflow and its representation evolved.

% The fourth sub-activity, \textit{justify and reflect}, involved reviewing the design decisions against the defined requirements. During development, features such as workflow-layer navigation, interactive filtering, summary calculation, and CSV export were assessed not only as interface functions, but also as responses to the identified lack of transparency, traceability, and usability in the original workflow. This reflection also helped identify possible directions for future development, including support for additional workflow stages and tighter integration with Volvo CE internal infrastructure.

% A key design decision in this thesis was to represent selected parts of the TMC workflow through a modular web-based prototype rather than attempt to replicate the SAP environment directly. Another important decision was to externalize selected workflow logic through backend services and SQL-based data handling, instead of relying on SAP-specific execution structures. This made the preparation, market view, adjustment, machine line split, and resplit-related logic easier to inspect and explain, while remaining manageable within the scope of the thesis.

% One challenge encountered during development was how to represent the complexity of the TMC workflow in a simplified but still meaningful way. The original SAP-based process involves multiple data sources, matrix-based business rules, and staged transformations. To make the artefact feasible within the thesis, the implementation was limited to workflow coverage from source data upload to machine line split, while preserving the essential logic needed to demonstrate the value of a more transparent and user-oriented representation.

% \paragraph{Alternative strategy}
% An alternative strategy would have been to develop the artefact through action research in closer integration with ongoing organizational processes. However, this was not chosen because the thesis was conducted in an exploratory mode and did not aim to introduce changes directly into the production environment. An iterative prototyping strategy combined with stakeholder feedback was therefore more suitable for developing and refining the artefact in this context.

% \subsection{Design and develop artifact}

% The third activity of the DSR framework is to design and develop the artefact, transforming the outlined artefact and elicited requirements into a concrete implementation. According to Johannesson and Perjons, this activity can be understood through four sub-activities: imagine and brainstorm, assess and select, sketch and build, and justify and reflect \cite{johannesson2021,hevner2004design}.

% The first sub-activity, \textit{imagine and brainstorm}, concerned how the TMC workflow could be represented in a more transparent and user-friendly way outside the SAP-based environment. Early ideas focused on data upload, workflow navigation, intermediate layer inspection, and calculation-related views.

% The second sub-activity, \textit{assess and select}, required evaluating alternative design directions. A web-based prototype was chosen over a script-based or spreadsheet-based solution, and a frontend-backend architecture was selected to separate user interaction from backend processing. The implementation scope was defined up to the machine line split stage.

% The third sub-activity, \textit{sketch and build}, translated these decisions into a working prototype. Uploaded CSV-based files were processed through backend services and stored in a local relational database, allowing selected workflow logic to be re-expressed through Python- and SQL-based processing in a more inspectable form than in the original SAP environment. The prototype was developed iteratively as the understanding of the workflow evolved.

% The fourth sub-activity, \textit{justify and reflect}, involved reviewing the design decisions against the defined requirements and checking whether the prototype produced results consistent with the intended workflow logic. Using 2024 TMC source data and matrix files, selected outputs were spot-checked against the defined calculation rules. This helped assess the correctness of the reformulated backend logic and identify remaining limitations of the artefact.




% A key design decision was to represent selected parts of the TMC workflow through a modular web-based prototype rather than attempt to replicate the SAP environment directly. Another was to externalize selected workflow logic through backend services and SQL-based data handling instead of relying on SAP-specific execution structures.

% One challenge during development was translating calculation logic embedded in the SAP environment into executable backend logic. In the original system, important workflow steps are expressed through planning functions, planning sequences, matrix-based structures, and staged transformations, none of which map directly onto reusable code. These structures had to be interpreted and reformulated into Python- and SQL-based processing logic. To keep the artefact feasible within the scope of the thesis, the implementation was limited to workflow coverage from source data upload to machine line split.

% \paragraph{Alternative strategy} An alternative strategy would have been to develop the artefact through action research in closer integration with ongoing organizational processes. However, given that the thesis was conducted in an exploratory mode without the aim of introducing changes into the production environment, iterative prototyping combined with stakeholder feedback was the more suitable approach.

\subsection{Design and Develop the Artefact}

The third DSR activity was to design and develop the artefact based on the requirements defined in the previous activity. Since the main contribution of this thesis is a working prototype, this activity formed a central part of the research process \cite{johannesson2021,hevner2004design}.

In this stage, CSV-based source files, matrix tables, mapping rules, workflow descriptions, and generated outputs were inspected to understand how selected parts of the workflow could be reconstructed outside the SAP environment. Stakeholder clarification was used when business rules, machine-line mappings, or data dependencies were unclear. The collected material was analysed through workflow decomposition, rule interpretation, and iterative output inspection, which supported the translation of selected workflow logic into explicit Python- and SQL-based backend processing steps.

The development followed an iterative prototyping approach. Early design ideas focused on staged outputs, workflow navigation, filtering, summary values, and export functions. These ideas were gradually implemented and refined as understanding of the workflow improved. The implemented scope covered source data upload, preparation, market view, adjustment, machine line split, and resplit. A key challenge was that important workflow logic in the original SAP-based process is embedded in planning functions, planning sequences, matrix-based structures, and staged transformations, so it had to be interpreted and reformulated into backend processing logic. The iterative approach was therefore not only a practical way of building software, but also a way of learning the workflow itself: each implemented stage made it easier to see what data, rule assumptions, and inspection needs had to be represented next.




% \subsection{Demonstrate artefact}





% The fourth activity of the DSR framework is to demonstrate the artefact. This activity shows how the artefact can be applied in a relevant case to address the previously explicated problem. In this sense, the demonstration serves as a proof of feasibility for the artefact \cite{johannesson2021,hevner2004design}.

% The first sub-activity was to choose a case on which to apply the artefact. In this thesis, the demonstration was based on a realistic TMC workflow scenario using 2024 TMC source data and matrix files. The case focused on how users could upload workflow-related input data, process it through the implemented stages, and inspect the resulting preparation-, adjustment-, and split-related outputs in the prototype.

% The second sub-activity was to apply the artefact to the selected case. The uploaded CSV-based source data and matrix files were processed through the backend, where selected workflow logic was handled through Python-based services and SQL-backed data processing. The resulting outputs were then made available in the web interface through layer-based views, interactive tables, summary calculations, and CSV export functionality.

% This demonstration showed that the artefact could be used to support the workflow from source data upload to machine line split in a more transparent and inspectable way than the corresponding SAP-based representation. It therefore served as a proof of concept for the feasibility of the proposed web-based prototype.

% \paragraph{Alternative strategy}
% An alternative strategy would have been a broader pilot deployment in an ongoing Volvo CE workflow. However, this was not chosen because the thesis focused on demonstrating the feasibility of the prototype in a realistic use case rather than introducing it into wider operational use.

% \subsection{Demonstrate the Artefact}

% The fourth DSR activity was to demonstrate the artefact. This activity shows how the artefact can be applied in a relevant case and can serve as a proof of feasibility for the proposed solution \cite{johannesson2021,hevner2004design}. In this thesis, the demonstration was based on a realistic TMC scenario using Volvo CE's 2024 market data and matrix files used for the yearly TMC calculation.

% The selected data collection methods were applied by using the prototype to upload source and matrix files, process them through the implemented backend stages, and generate outputs for preparation, market view, adjustment, machine line split, and resplit. These outputs were then made available through the web interface using layer-based views, interactive tables, summary values, charts, and CSV export functions.

% The demonstration focused on showing how users could follow selected workflow stages step by step, inspect intermediate outputs, and review results through the prototype interface. In this way, the activity demonstrated the feasibility of representing selected parts of the TMC workflow outside the SAP environment in a more transparent and inspectable form.

% The detailed demonstration of the prototype and its generated workflow outputs is presented in the results chapter.


% \subsection{Demonstrate the Artefact}

% The fourth DSR activity was to demonstrate the artefact. This activity shows how the artefact can be applied in a relevant case and can serve as a proof of feasibility for the proposed solution \cite{johannesson2021,hevner2004design}. In this thesis, the demonstration was based on a realistic TMC scenario using Volvo CE's 2024 market data and matrix files used for the yearly TMC calculation.

% The selected data collection methods were applied by using the prototype to upload source and matrix files, process them through the implemented backend stages, and generate preparation, market view, adjustment, machine line split, and resplit outputs. These outputs were then made available through the web interface using layer-based views, interactive tables, summary values, and CSV export functions.

% The generated outputs were analysed through structured output inspection and comparison with the expected workflow logic. In particular, the results were checked against the defined business rules, expected intermediate outputs, and, where comparable, the results produced by the existing SAP-based TMC implementation. This helped assess whether the reformulated backend logic behaved consistently with the intended TMC process within the implemented scope.

% This activity served as a proof of feasibility for the artefact within the selected scope, from source data upload to machine line split. The detailed demonstration, output inspection, and comparison results are presented in the results and evaluation chapters.


\subsection{Demonstrate the Artefact}

The fourth DSR activity was to demonstrate the artefact. This activity shows how the artefact can be applied in a relevant case and can serve as a proof of feasibility for the proposed solution \cite{johannesson2021,hevner2004design}. In this thesis, the artefact was demonstrated through the Volvo CE case study using a realistic TMC scenario with 2024 market data and matrix files.

The prototype was used to upload source and matrix files, process them through the implemented backend stages, and generate outputs for preparation, market view, adjustment, machine line split, and resplit processing. These outputs were made available through the web interface using layer-based views, interactive tables, summary values, charts, and CSV export functions. The activity therefore demonstrated the feasibility of representing selected parts of the workflow outside the SAP environment in a more transparent and inspectable form.

The detailed demonstration of the prototype and its generated workflow outputs is presented in the results chapter.



% \subsection{Evaluate Artefact}
% The final activity of the DSR framework is to evaluate the artefact. This activity aims to determine whether the artefact fulfils its intended purpose and satisfies the previously defined requirements.\cite{johannesson2021} In this thesis, the artefact was evaluated through prototype use, output inspection, and stakeholder feedback.

% The first sub-activity was to analyse the evaluation context. The evaluation was carried out using the deployed web-based prototype, which was made accessible to stakeholders through a test deployment. Representative 2024 TMC source data and matrix files were used as input. The setting was realistic in terms of data and workflow context, but still limited to the current prototype scope.

% The second sub-activity was to select evaluation goals and strategy. The main goals were to assess whether the artefact improved transparency, usability, and inspectability compared with the SAP-based TMC implementation, and whether it fulfilled the defined requirements. Since the artefact had not yet been introduced into broader internal use, the evaluation was formative and qualitative in nature.

% The third sub-activity was to design and carry out the evaluation. A feedback session was conducted with a small group of two to three stakeholders from Volvo CE, including personnel from the Strategy and Business Intelligence team responsible for market research and the yearly TMC cycle. Stakeholders accessed the deployed prototype directly through a test environment and interacted with the system using real workflow data. Feedback was collected through semi-structured discussion following the session.

% The evaluation indicated that the artefact improved workflow visibility and made intermediate results easier to inspect compared with the SAP-based process. Stakeholders found the web-based interface easier to operate than the existing SAP environment and noted that the layer-based views provided clearer visibility into intermediate workflow steps than was previously available. The interactive filtering, summary calculation, and CSV export features were considered useful for review and validation purposes. Stakeholders also highlighted the potential for further development, in particular the extension of the prototype to cover later workflow stages such as Total Market Calculation, suggesting that the current implementation provides a solid foundation for broader future use.

% \paragraph{Alternative strategy}
% An alternative strategy would have been a broader pilot deployment within an ongoing Volvo CE workflow. However, this was not chosen because the artefact is still at the prototype stage, and a formative qualitative evaluation was more suitable for identifying improvements before wider internal use.



% \subsection{Evaluate the Artefact}

% The final DSR activity was to evaluate the artefact. This activity aimed to determine whether the prototype fulfilled its intended purpose and satisfied the requirements defined earlier in the research process \cite{johannesson2021,hevner2004design}. Since the artefact was developed as a research prototype, the evaluation was formative and qualitative rather than a full production assessment.

% In this activity, data were collected through prototype use, structured output inspection, and stakeholder feedback. The prototype was evaluated in a test setting using Volvo CE's 2024 market data and matrix files. The generated outputs were inspected to assess whether the implemented workflow stages behaved according to the defined business rules and expected workflow logic. Where comparable, selected results were also checked against the existing SAP-based TMC process.

% The collected material was analysed through requirement-based evaluation and qualitative feedback analysis. The evaluation focused on whether the artefact supported transparency, traceability, usability, validation support, and prototype-level performance within the implemented scope. Stakeholder feedback was used to understand whether the prototype made intermediate workflow stages clearer and easier to review from a user perspective.

% The detailed evaluation design, parameters, results, and limitations are presented in the evaluation chapter.



\subsection{Evaluate the Artefact}

The final DSR activity was to evaluate the artefact. This activity aimed to determine whether the prototype fulfilled its intended purpose and satisfied the requirements defined earlier in the research process \cite{johannesson2021,hevner2004design}. Since the artefact was developed as a research prototype, the evaluation was formative and qualitative rather than a full production-level assessment.

Data were collected through prototype use, structured output inspection, and stakeholder feedback. The prototype was evaluated in a test setting using Volvo CE's 2024 market data and matrix files, and the evaluation focused on the implemented workflow scope from source data upload to preparation, market view, adjustment, machine line split, and resplit.

The generated outputs were inspected to assess whether the implemented stages behaved according to the defined business rules and expected workflow logic. Where comparable, selected results were also checked against outputs from the existing SAP-based TMC process. The collected material was analysed through requirement-based evaluation and qualitative feedback analysis, with attention to functional fulfilment, prototype-level performance, transparency, traceability, usability, and validation.

The detailed evaluation parameters, results, and limitations are presented in the evaluation chapter.











\section{Ethical Considerations}

This study was conducted in collaboration with Volvo CE and required careful attention to research integrity, confidentiality, and responsible handling of research material. The study was guided by principles of good research practice, including transparency, honesty, and respect in the conduct and reporting of research \cite{su_ethics_policy}.

The empirical part of the thesis involved workflow analysis, document review, prototype development, and recurring review meetings and clarification discussions with relevant stakeholders. No sensitive personal data was collected or processed, and the study did not involve any physical intervention or any attempt to influence participants physically or mentally. Stakeholders who contributed contextual information were informed about the purpose of the study, and participation in discussions was voluntary \cite{etikprovningsmyndigheten_act}.

The purpose of these interactions was solely to understand current work practices and support the design and assessment of the artefact, not to evaluate or report individual performance or behaviour. In this way, the study sought to respect participants' autonomy while limiting the use of collected material to what was necessary for the research purpose \cite{vr_good_research_2024}.

Special consideration was also given to the handling of company-internal information. The prototype was developed and demonstrated in a controlled setting and was not used as a production replacement. Where screenshots, figures, code fragments, system outputs, or data samples from internal systems were included for explanatory purposes, they were anonymised or masked to remove company-specific identifiers and potentially sensitive information. The thesis does not disclose Volvo CE trade secrets, confidential business data, or complete proprietary calculation logic, and the handling of company-related material followed Volvo CE's internal requirements for confidentiality and responsible information handling \cite{su_ethics_policy}.

Overall, the study was conducted in accordance with Stockholm University's ethical framework, Swedish principles for good research practice, and Volvo CE's internal requirements for confidentiality and responsible information handling, with particular attention to confidentiality, responsible data handling, and respect for the individuals involved in the research context \cite{vr_good_research_2024}.


























